\section{Cusp jets on the covers used in the degree estimate}
\label{sec:cusp-jets}

The geometric estimate takes place on a fixed good cover in a family.
Fix a sufficiently small positive \(\epsilon\), take a family meeting
the largeness conditions for this \(\epsilon\), and let \(x_1\) be its
lower endpoint. There exists \(x_0\) such that the following statements
hold for every \(x>x_1\) with \(x\ge x_0\), every index \(\iota\)
above \(x\), every point satisfying the prescribed point and
reciprocity conditions, and every stage \(i\) with \(i+1<m\).
For each such stage they hold for every good cover \(Y\).
The bounds for \(x_0\) can depend on the family. In the conjunction
of two estimates the upper bound on \(\epsilon\) is the minimum of
their admissible upper bounds, the largeness bounds are their
componentwise maxima, and the eventual threshold is the maximum of
the two \(x_0\)'s. Thus one can apply all the estimates below at the
same \(x,\iota,i,Y\). In a later restriction to an arithmetic
window, the restriction must lie above this eventual threshold.
\footnote{\href{https://github.com/Deicyde/riemann-hypothesis/blob/40e56bab/src/Development/AutomorphicArithmetic/ThetaCorrespondence/ThetaPairLifts/ThirdPlaceSeries/LiftLambdaBound.lean\#L17-L32}{Quantifiers of the geometric cover condition},
\href{https://github.com/Deicyde/riemann-hypothesis/blob/40e56bab/src/Development/AutomorphicArithmetic/ThetaCorrespondence/ThetaPairLifts/ThirdPlaceSeries/ExistsCuspWithThetaJetValueLe.lean\#L24-L43}{simultaneous admissible range and eventual threshold}.}

\subsection{A theta rule of bounded valuation}

Let \(E\) be the CM host field of degree \(d=[E:\mathbb Q]\).
For a dominant array \(a\) of theta exponents on the datum, write
\(\delta=D_F(a)\) for its nonnegative degree as used in the jet
estimate. Let \(n\ge1\) be the number of theta columns, at least
the sum of twice the ranks of the indicated lattice summands. Put
\[
            R=\frac{|\operatorname{disc}E|^{2/d}}{2d}.
\]
On the characteristic-zero-residue integral chart of the proof,
let \(T\) be the complete discrete valuation ring, \(v_T\) its
integer-valued order with \(v_T(0)=+\infty\), and \(\phi_T\ge0\) its
given slope parameter. The chart has a truncated frame and a PEL
uniformization with the specified torsion isomorphism. The auxiliary
element \(\alpha\in\mathcal O_E\) satisfies
\(\alpha\overline\alpha=q\), with \(q\) dividing the level and the
level ideal contained in \((\overline\alpha)\).

Let \(N\) bound \(d\) times the negative signature rank of the
datum. Suppose the numerical data satisfy
\begin{equation}\label{eq:jets-chart-num}
    4\delta+4\le q,\qquad 2\delta\sqrt N<q,\qquad
    \frac{\delta^2 R N^3}{q}\le c .
\end{equation}
Then some theta rule \(f\) of type \((a,n)\) has nonzero evaluation
on this frame, with
\begin{equation}\label{eq:jets-value}
                   v_T(f)\le c\,\phi_T .
\end{equation}

The inputs to this inference are the characteristic-zero harmonic
polynomial with local invariant at most zero, the existence of the
power data and short basis on the chart, and the local comparison of
the infimum of the normalized theta valuations with the local
invariant. Fix a uniformizer and such a polynomial. The short
basis has parameter \(0\le\rho\le N\), so that
\(2\delta\sqrt{\rho}<q\). Apply the local comparison with valuation
normalization \(1\) and with these choices. The local invariant term
is nonpositive, leaving
\[
    \inf_f v_T(f)\le B,\qquad
    B=\frac{\delta^2 R\rho^3\phi_T}{q}
       \le \frac{\delta^2 RN^3}{q}\phi_T\le c\phi_T .
\]
The infimum runs over the rules obtained from this polynomial.
If every one of them had zero evaluation or order strictly larger
than \(c\phi_T\), each extended order would be at least
\(\lfloor B\rfloor+1\). The same would hold for the infimum,
contradicting \(\lfloor B\rfloor+1>B\). Thus a member with nonzero
evaluation satisfies \eqref{eq:jets-value}. The discreteness at this
step avoids any loss in passing from an infimum to a section.
\footnote{\href{https://github.com/Deicyde/riemann-hypothesis/blob/40e56bab/src/Development/AbelianVarietyTheory/Proofs/DivisorProofs/ReflectionSumJetNumber.lean\#L52-L98}{The short-basis, harmonic-polynomial and infimum inputs, followed by the argument with the integral valuation}.}

In the family application \(q=g_0g_1\), the product of the two host
primes, and \(N=2fm\) in the rank notation of the family. The
numerical jet condition at \(g_0\) supplies the first two inequalities
of \eqref{eq:jets-chart-num} with denominator bound \(g_0\), together
with
\[
       \frac{RN^3}{g_0}\left(\frac{rN}{4}+2\delta^2\right)\le c.
\]
Here \(r\ge0\) is the rank parameter in that condition. As
\(q\ge g_0>0\) and \(RN^3\ge0\),
\[
       \frac{\delta^2RN^3}{q}
       \le\frac{\delta^2 RN^3}{g_0}
       \le\frac{RN^3}{g_0}\left(\frac{rN}{4}+2\delta^2\right).
\]
The first two inequalities remain valid on replacing \(g_0\) by
\(q\). This proves the three bounds required on the chart.
\footnote{\href{https://github.com/Deicyde/riemann-hypothesis/blob/40e56bab/src/Development/AbelianVarietyTheory/Proofs/DivisorProofs/ReflectionSumJetNumber.lean\#L42-L50}{The degree comparison},
\href{https://github.com/Deicyde/riemann-hypothesis/blob/40e56bab/src/Development/AutomorphicArithmetic/ThetaCorrespondence/ThetaPairLifts/ThirdPlaceSeries/ExistsThetaRuleValueLeCoverCharZero.lean\#L47-L72}{specialization to the two host primes and the field degree}.}

\subsection{Descent to the cusp and the factor \(1/2500\)}

Take a component \(Z\), a toroidal flag chart, a boundary point
\(y\) outside its distinguished open locus, and a local map of the
stalk at \(y\) to a complete discrete valuation ring, compatible
with the given generic frame. The local map on this chart implies
that the residue field has characteristic zero. The chart-extension
theorem provides a valuation ring \(T'\), an extended frame admitting
the truncated uniformizing chart, and a positive integer \(e_T\).
The slope of the extended frame is \(e_T\phi_T\), and its extended
order on evaluations of a theta rule is \(e_T\) times the order on
the original frame. Use the canonical connection in this comparison.

On the extended frame choose \(f\) by \eqref{eq:jets-value}.
Its value there is nonzero and its order is at most
\(c e_T\phi_T\). If its value on the original frame were zero, the
comparison of extended orders would identify a finite number with
\(+\infty\). Thus its order there is finite as well. Divide the
order comparison by \(e_T>0\) to get the original bound \(c\phi_T\).
By dominance and positivity of the transformed theta weight, the
theta cuspidality theorem makes this a cusp form. Restriction to
the component gives a cusp form with the same evaluation because
the chosen point factors through \(Z\).
\footnote{\href{https://github.com/Deicyde/riemann-hypothesis/blob/40e56bab/src/Development/AutomorphicArithmetic/ThetaCorrespondence/ThetaPairLifts/ThirdPlaceSeries/ExistsCuspWithThetaJetValueLe.lean\#L44-L65}{Descent of nonvanishing and the valuation bound},
\href{https://github.com/Deicyde/riemann-hypothesis/blob/40e56bab/src/Development/AutomorphicArithmetic/ThetaCorrespondence/ThetaPairLifts/ThirdPlaceSeries/ThetaValueOnCuspNeighbourhood.lean\#L29-L57}{cuspidality and the equality after restriction to the component}.}

For the two weights in use, denote the weight of the quotient by
\(u_1\) and its companion upper weight by \(u_b\). The theta
construction first shifts each weight by its half-column weight.
There are \(n\) columns with \(8\mid n\), so \(4\mid n/2\).
The four-divisibility boundary comparison removes this shift and
preserves the nonzero value's order. The numerical estimate for
\(u_1\) follows from the family size inequalities. The upper-weight
degree comparison gives a theta degree for \(u_b\) no larger than
the one for \(u_1\), so the same numerical jet bound works for both.
Its coefficient is \(c=1/2500\). Extending the cusp form to the
compactified flag space under the extension isomorphism preserves its
evaluation. Thus for each of these two weights there is a section
\(w\) in the specified cusp subspace with
\[
            w(y')\ne0,\qquad v_T(w(y'))\le\phi_T/2500.
\]
Only the section is chosen after the component, boundary point and
valuative frame.
\footnote{\href{https://github.com/Deicyde/riemann-hypothesis/blob/40e56bab/src/Development/AutomorphicArithmetic/ThetaCorrespondence/ThetaPairLifts/ThirdPlaceSeries/JetsOfConstantSection.lean\#L26-L49}{The shift by the half-column weight and the boundary comparison},
\href{https://github.com/Deicyde/riemann-hypothesis/blob/40e56bab/src/Development/AutomorphicArithmetic/ThetaCorrespondence/ThetaPairLifts/ThirdPlaceSeries/JetNumUb.lean\#L23-L34}{comparison of the two degrees},
\href{https://github.com/Deicyde/riemann-hypothesis/blob/40e56bab/src/Development/AutomorphicArithmetic/ThetaCorrespondence/ThetaPairLifts/ThirdPlaceSeries/ExistsCuspValueLePhiOfHilbertSamuel.lean\#L24-L50}{the two weights, simultaneous threshold and final valuation}.}

\subsection{From boundary orders to the degree comparison}
Let \(\nu\) be a discrete valuation centered at a boundary point. Its
potential \(\phi\) is represented by a valuative frame and a real \(h>0\)
for which \(\phi=h\phi_T\). The additive order of a germ of nonzero
value is \(h\) times the order of its pullback to the valuation ring.
Choose the cusp section constructed above. A change of local
trivialization multiplies its representative by a unit, so
\[
  \operatorname{ord}_{\nu}(w)=h v_T(w(y'))\le h\phi_T/2500=\phi/2500 .
\]
An independent Hodge-section theorem gives a nonzero section \(\theta\)
of weight \(8(p+q)\) times the Hodge weight whose order is at least
\(8\phi/625\) wherever its germ has finite order. Here \(p,q\) are the
two signature ranks. Since \(\phi\ge0\), we obtain
\(3\operatorname{ord}_\nu(\theta)\ge8\operatorname{ord}_\nu(w)\).
In the multiplicative convention this gives
\(\nu(\theta)^3\le\nu(w)^8\). If \(\nu(\theta)=0\), the same inequality
is immediate. Combining this with base-point freeness on the open
locus and extension of cusp forms supplies the valuation-multiplicity
premise in the component-image degree comparison. That comparison
also uses compactification proportionality, the top Chern number
formula and birationality at the prescribed weight gaps.
\footnote{\href{https://github.com/Deicyde/riemann-hypothesis/blob/40e56bab/src/Development/AutomorphicArithmetic/ThetaCorrespondence/ThetaPairLifts/ThirdPlaceSeries/ExistsSectionOrderBoundHilbertSamuel.lean\#L22-L46}{Passing to germs};
\href{https://github.com/Deicyde/riemann-hypothesis/blob/40e56bab/src/Development/AutomorphicArithmetic/ThetaCorrespondence/ThetaPairLifts/ThirdPlaceSeries/ExistsNonzeroSectionWeightBound.lean\#L79-L82}{the Hodge section};
\href{https://github.com/Deicyde/riemann-hypothesis/blob/40e56bab/src/Development/AutomorphicArithmetic/ThetaCorrespondence/ThetaPairLifts/ThirdPlaceSeries/ValuationMultiplicityOfConsecutiveRun.lean\#L48-L97}{the powers three and eight};
\href{https://github.com/Deicyde/riemann-hypothesis/blob/40e56bab/src/Development/AutomorphicArithmetic/ThetaCorrespondence/ThetaPairLifts/ThirdPlaceSeries/CoveringDegreeMultiplicative.lean\#L25-L56}{application to the cover degree}.}

\subsection{The metric loss in the lift bound}
Write \(X,b\) for the size parameters with \(b=\exp L\). The Gram
condition identifies the squared norm on the integral cusp lattice with
the norm of the chosen frame. The Tamagawa measure comparison bounds
the maximal slope by
\[
 2\lambda_{\rm sup}+2\log C_{\rm sup}+
 [F:\mathbb Q]^{-1}\log|\operatorname{disc}F|,
\]
where \(F\) is the coefficient field. The three inputs are the mean
logarithmic sup-norm estimate, the logarithm of the supremum constant,
and the root-discriminant estimate. Let \(A\) be the size exponent.
The weight parameters obey \(0\le s\le8frb\) and \(0\le t\le b\).
Since \(fr,b\ge1\), this gives \(s+t+1\le10frb\). The logarithmic
estimates on these parameters and \(frL\le X^A\le b\) enter the
quantitative comparison for the three terms above, which bounds their
sum by \(X^4(X^A)^2b=X^{2A+4}b\). Thus the cusp-slope exponent is
\(A_1=2A+4\). The all-places measure comparison of constant \(C_m\)
and the host estimate of exponent \(A_h\) give the lift bound with
exponent \(A_1+(A_h+1)+\lceil2C_m\rceil\). Taking the maximum of
the cusp, measure, host, Gram and weight thresholds makes all five
inputs applicable at the same point.
\footnote{\href{https://github.com/Deicyde/riemann-hypothesis/blob/40e56bab/src/Development/AutomorphicArithmetic/ThetaCorrespondence/FockModels/TamagawaPeterssonLambdaSqNonneg.lean\#L596-L669}{Metric comparison};
\href{https://github.com/Deicyde/riemann-hypothesis/blob/40e56bab/src/Development/AutomorphicArithmetic/ThetaCorrespondence/ThetaPairLifts/ThirdPlaceSeries/CoveringLambdaInvariantBound.lean\#L21-L37}{all-places exponent};
\href{https://github.com/Deicyde/riemann-hypothesis/blob/40e56bab/src/Development/AutomorphicArithmetic/ThetaCorrespondence/ThetaPairLifts/ThirdPlaceSeries/TopCovolG.lean\#L50-L58}{degree and metric bounds used in the covolume argument}.}
